% ============================================================
%  ch04 —— §2.3（后半）–2.5（原文 p.16–p.18）
% ============================================================
\routenote{p.16--p.18}

几何 Hecke 算子定义为函子（见 [17]）：
\begin{align*}
\Hk:\ \Rep(\ckG)\times D^b(\Bun_{\mathbb{G}(m,n)})&\longrightarrow D^b\big(\Bun_{\mathbb{G}(m,n)}\times\PP^1\setminus\{0,\infty\}\big),\\
(V,K)&\longmapsto \Hk_V(K):=\mathrm{pr}_{2,!}\mathrm{pr}_1^*(K\otimes\mathrm{IC}_V).
\end{align*}
为了复合这些算子，把 $\Hk_V$ 扩充为算子 $D^b(\Bun_{\mathbb{G}(m,n)}\times\PP^1\setminus\{0,\infty\})\to D^b(\Bun_{\mathbb{G}(m,n)}\times\PP^1\setminus\{0,\infty\})$：
\[
K\longmapsto \mathrm{pr}_{2,!}\Big(\big(\mathrm{pr}_1\times\mathrm{pr}_{\PP^1\setminus\{0,\infty\}}\big)^*K\otimes\mathrm{IC}_V\Big).
\]

设 $E$ 是 $\PP^1\setminus\{0,\infty\}$ 上的 $\ckG$-局部系统，视作张量函子
\[
E:\ \mathcal S\cong\Rep(\ckG)\to\Loc(\PP^1\setminus\{0,\infty\}),
\]
把 $V\in\Rep(\ckG)$ 送到 $E_V\in\Loc(\PP^1\setminus\{0,\infty\})$。

以 $E$ 为特征值的 Hecke 特征层是 perverse 层 $K\in\mathrm{Perv}(\Bun_{\mathbb{G}(m,n)})$ 连同同构
\[
\Hk_V(K)\xrightarrow{\ \sim\ }K\boxtimes E_V,
\]
这些同构与 $\Rep(\ckG)$ 上的对称张量结构以及 Hecke 对应的复合相容（关于相容性公理的详细阐述见 [17]）。

\begin{remark}
由于局部系统通常是以别的方式引入的，让我们简要回忆张量函子 $E$ 如何重构局部系统的单值表示。这对于固定单值表示的记号是有用的。取 $\PP^1\setminus\{0,\infty\}$ 中 $\operatorname{Spec}K_0$ 上方的一个几何点 $\eta$。限制到 $\eta$ 定义张量函子
\begin{equation}
\omega_E:\ \mathcal S\xrightarrow{\ E\ }\Loc(\PP^1\setminus\{0,\infty\})\xrightarrow{\ j_\eta^*\ }\mathrm{Vec},
\end{equation}
亦即（见 [12, Theorem 3.2]）一个带 $\pi_1(\PP^1\setminus\{0,\infty\},\eta)$ 作用的 $\ckG$-挠子。取该挠子的一个点，便定义了单值表示
\begin{equation}
\phi:\ \pi_1\big(\PP^1\setminus\{0,\infty\},\eta\big)\to\ckG(\mathbb{Q}_\ell).
\end{equation}
记 $\phi^{\mathrm{geo}}$ 为 $\phi$ 到 $\pi_1^{\mathrm{geo}}(\PP^1\setminus\{0,\infty\},\eta)=\pi_1(\PP^1\setminus\{0,\infty\}\otimes_k\bar k,\eta)$ 的限制，称之为 $E$ 的整体几何单值表示。
\end{remark}

上述构造对扭群 $\mathbb{G}$ 的类比将产生 $\LG$-局部系统。这里把 $\mathbb{G}$ 的 $L$-群定义为 $\LG=\ckG\rtimes\langle\sigma\rangle$，其中 $\sigma$ 是 $\mathbb{G}$ 的定义中所用的 $G$ 的 $N$ 阶自同构。这个自同构通过几何 Satake 同构诱导 $\ckG$ 的自同构。在引理 9.3 中我们将验证这个自同构确实保持 $\ckG$ 的钉住。

让我们给出对本文目的足够的 $\LG$-局部系统的定义。我们不能立即应用几何 Satake 同构的原因是 $\mathrm{pr}_2$ 的纤维沿 $\PP^1\setminus\{0,\infty\}$ 不是常值的，故只有部分 Hecke 算子 $\Hk_V$ 会定义 $\PP^1\setminus\{0,\infty\}$ 上的整体 Hecke 算子。不过，我们可以沿映射 $[N]:\PP^1\setminus\{0,\infty\}\to\PP^1\setminus\{0,\infty\}$ 拉回卷积图。把这个覆盖记作 $\widetilde{\PP^1\setminus\{0,\infty\}}\to\PP^1\setminus\{0,\infty\}$。

拉回之后，我们可以像以前一样在 $\Hecke^{\widetilde{\PP^1\setminus\{0,\infty\}}}_{\mathbb{G}(m,n)}$ 上定义 Hecke 算子
\[
\Hk:\ \Rep(\ckG)\times D^b(\Bun_{\mathbb{G}(m,n)})\cong\mathcal S\times D^b(\Bun_{\mathbb{G}(m,n)})\to D^b\big(\Bun_{\mathbb{G}(m,n)}\times\widetilde{\PP^1\setminus\{0,\infty\}}\big).
\]
此外，覆盖群 $\mu_N$ 作用在 $\widetilde{\PP^1\setminus\{0,\infty\}}$ 上的卷积图上。这在函子 $\Hk$ 上定义了 $\mu_N$-等变结构：在源上，$\mu_N$ 通过 $\sigma:\mu_N\to\operatorname{Aut}^\dagger(\mathbb{G})$ 作用在 $\mathcal S$ 上，而后者可等同于 $\operatorname{Aut}^\dagger(\ckG)$ 在 $\Rep(\ckG)$ 上的作用（引理 9.3）；在靶上，$\mu_N$ 作用在 $\widetilde{\PP^1\setminus\{0,\infty\}}$ 上。

设 $E$ 是 $\widetilde{\PP^1\setminus\{0,\infty\}}$ 上的 $\ckG$-局部系统，连同对 $\zeta\in\mu_N$ 的相容同构 $\zeta^*E\cong E\times^{\ckG,\sigma(\zeta)}\ckG$。我们把 $E$ 视作带 $\mu_N$-等变结构的张量函子
\[
E:\ \Rep(\ckG)\to\Loc(\widetilde{\PP^1\setminus\{0,\infty\}}).
\]
和前一样，可以定义以 $E$ 为特征值的 Hecke 特征层 $K\in\mathrm{Perv}(\Bun_{\mathbb{G}(m,n)})$，但现在须指定与 $\Rep(\ckG)$ 上张量结构相容的函子同构
\[
\epsilon(V):\ \Hk_V(K)\xrightarrow{\ \sim\ }K\boxtimes E_V,
\]
它与函子 $V\mapsto\Hk_V(K)$ 和函子 $E$ 的 $\mu_N$-等变结构交换。

注意，若 $\sigma:\mu_N\to\operatorname{Aut}(\mathbb{G})$ 平凡，则同构 $\zeta^*E\cong E\times^{\ckG,\sigma(\zeta)}\ckG$ 给出 $E$ 的下降数据。故此时定义与常群情形一致。

有了这些定义，我们可以陈述第一个主要结果：

\begin{mainthm}[1]
\begin{enumerate}[label=(\arabic*), leftmargin=1.9em, itemsep=.25em]
\item 层 $A=A_\varphi$ 与 $A_{\varphi,\chi}$ 是 Hecke 特征层。记 $A_\varphi$（相应地 $A_{\varphi,\chi}$）的特征值为 $\Kl^{\LG}(\varphi)$（相应地 $\Kl^{\LG}(\varphi,\chi)$）。
\item 若 $\mathbb{G}=\PP^1\times G$ 是常群概形，则局部系统 $\Kl^{\ckG}(\varphi)$ 在 $0$ 处驯顺分歧。$\Kl^{\ckG}(\varphi)$ 在 $0$ 处的单值作用由 $\ckG$ 中的一个主幂幺元给出。
\item 对任意不可约表示 $V\in\Rep(\ckG)$，层 $\Kl^{\LG}(\varphi,\chi)^V$ 是纯的。
\end{enumerate}
\end{mainthm}

由于我们用几何 Hecke 算子定义了 $\Kl^{\ckG}(\varphi)$，对任意点 $x\in\PP^1\setminus\{0,\infty\}$，由该局部系统定义的 $\mathrm{Frob}_x$ 的 $\ckG$-共轭类由 Gross 的自守形式 $\pi(\varphi)$ 在 $x$ 处的 Satake 参数给出。

\subsection{单值表示}

本节假设 $\mathbb{G}=G\times\PP^1$，其中 $G$ 是 $k$ 上几乎单的分裂群。如注 2.10 那样，记 $\phi$ 为 $\Kl^{\ckG}(\varphi)$ 的单值表示，$\phi^{\mathrm{geo}}$ 为其几何单值表示。

下一个结果是关于 $\Kl^{\ckG}(\varphi,\chi)$ 在 $\infty$ 处的几何单值的。记 $\Kl^{\mathrm{Ad}}_{\ckG}(\varphi,\chi)$ 为由伴随表示 $V=\ckg$ 诱导的 $\PP^1\setminus\{0,\infty\}$ 上的局部系统。固定局部 Galois 群 $\operatorname{Gal}(K^{\mathrm{sep}}_\infty/K_\infty)$ 到 $\pi_1(\PP^1\setminus\{0,\infty\},\eta)$ 的一个嵌入，便得到 $\operatorname{Gal}(K^{\mathrm{sep}}_\infty/K_\infty)$ 在 $\Kl^{\mathrm{Ad}}_{\ckG}(\varphi,\chi)$ 于形式去心圆盘 $\operatorname{Spec}K^{\mathrm{sep}}_\infty$ 处的几何茎上的作用。把这个茎与 $\ckg$ 等同，便得到 $\operatorname{Gal}(K^{\mathrm{sep}}_\infty/K_\infty)$ 在 $\ckg$ 上的作用（在 $\ckG$-共轭意义下良定义）。

设 $I_\infty\subset\operatorname{Gal}(K^{\mathrm{sep}}_\infty/K_\infty)$ 为惯性群，$I^+_\infty\subset I_\infty$ 为野惯性群，$I^t_\infty=I_\infty/I^+_\infty$ 为驯惯性群。

\begin{mainthm}[2]
设 $p=\operatorname{char}(k)$ 对 $G$ 良好（即当 $G$ 是 $B_n$、$C_n$ 型时 $p>2$，当 $G$ 是 $F_4$、$G_2$ 型时 $p>3$）。则
\begin{itemize}[leftmargin=1.6em, itemsep=.2em]
\item $\Swan_\infty\big(\Kl^{\mathrm{Ad}}_{\ckG}(\varphi,\chi)\big)=r(\ckG)$，即 $\ckG$ 的秩；
\item $\ckg^{I_\infty}=0$。
\end{itemize}
\end{mainthm}

Swan 等式将在推论 5.1 中证明；$I_\infty$-不变量的消没将在命题 5.3(2) 中证明。

\begin{corollary}
（Gross--Reeder [20, Proposition 5.6]）设 $p=\operatorname{char}(k)$ 不整除 $\#W$。则局部几何 Galois 表示 $\phi^{\mathrm{geo}}_\infty:I_\infty\to\ckG$ 是 [20, §6] 中定义的简单野参数。更确切地说，在 $\ckG$-共轭意义下，有正合列的交换图
\[
\begin{tikzcd}[column sep=large, row sep=small]
I^+_\infty \arrow[r] \arrow[d, "\phi^{\mathrm{geo},+}_\infty"'] & I_\infty \arrow[r] \arrow[d, "\phi^{\mathrm{geo}}_\infty"] & I^t_\infty \arrow[d] \\
\ckT \arrow[r] & N(\ckT) \arrow[r] & W
\end{tikzcd}
\]
其中
\begin{itemize}[leftmargin=1.6em, itemsep=.2em]
\item 驯惯性群 $I^t_\infty$ 的一个拓扑生成元映到 Coxeter 元 $\mathrm{Cox}\in W$（在 $W$-共轭意义下良定义，见 [6, Ch.~V, §6]）；
\item 野惯性群 $I^+_\infty$ 满射到一个子群 $\ckT(\zeta)\subset\ckT[p]$。这里 $\zeta\in\overline{\mathbb{F}}_p^\times$ 是一个本原 $h$ 次单位根（$h$ 为 Coxeter 数），而 $\ckT(\zeta)\subset\ckT[p]$ 是 $\ckT[p]$ 中唯一的 $\mathbb{F}_p[\mathrm{Cox}]$-子模，它同构于 $\mathbb{F}_p[\zeta]$（$\mathrm{Cox}$ 在其上以 $\zeta$ 乘法作用）。
\item $I_\infty$-表示 $\Kl^{\mathrm{Ad}}_{\ckG}(\varphi,\chi)$ 的非零跳跃等于 $1/h$。
\end{itemize}
\end{corollary}

最后我们可以陈述关于 $\Kl^{\ckG}(\varphi)$ 的整体几何单值的结果。设 $\ckG^{\mathrm{geo}}\subset\ckG$ 为整体几何单值表示 $\phi^{\mathrm{geo}}$ 之像的 Zariski 闭包。

\begin{mainthm}[3]
设 $\operatorname{char}(k)>2$。则 $\Kl^{\ckG}(\varphi)$ 的几何单值群 $\ckG^{\mathrm{geo}}$ 连通，且
\begin{itemize}[leftmargin=1.6em, itemsep=.2em]
\item 若 $\ckG$ 不是 $A_{2n}$（$n\ge 2$）型或 $B_3$ 型，则 $\ckG^{\mathrm{geo}}=\ckG^{\operatorname{Aut}^\dagger(\ckG),\circ}$；
\item 若 $\ckG$ 是 $A_{2n}$ 型，则 $\ckG^{\mathrm{geo}}=\ckG$；
\item 若 $\ckG$ 是 $B_3$ 型且 $\operatorname{char}(k)>3$，则 $\ckG^{\mathrm{geo}}=G_2$。
\end{itemize}
\end{mainthm}

证明将在 §6.2 给出。

\subsection{变体}

我们的构造有一个用 $D$-模代替 $\ell$-进层的变体。此时基域取为 $k=\mathbb{C}$。Artin--Schreier 局部系统 $\mathrm{AS}_\psi$ 被 $\Aff^1_\mathbb{C}=\operatorname{Spec}\mathbb{C}[x]$ 上的指数 $D$-模 $\mathbb{C}\langle x,\partial_x\rangle/(\partial_x-1)$ 替代，构造的其余部分照搬，于是得到张量函子
\[
\Kl^{\LG}(\varphi)_{dR}:\ \Rep({}^L G_\mathbb{C})\to\mathrm{Conn}(\PP^1\setminus\{0,\infty\},\mathbb{C}),
\]
其中 $\mathrm{Conn}(\PP^1\setminus\{0,\infty\},\mathbb{C})$ 是 $\PP^1\setminus\{0,\infty\},\mathbb{C}$ 上带联络的向量丛的张量范畴。

我们猜想本文的构造应给出与 Frenkel--Gross 构造相同的联络。为精确陈述，设 $G$ 几乎单，$\mathbb{G}$ 是 $G$ 在 $\PP^1\setminus\{0,\infty\}$ 上的拟分裂形式，由 $\sigma:\mu_N\to\operatorname{Aut}^\dagger(G)$ 给出。回忆 [15, §5] 中 $\PP^1\setminus\{0,\infty\},\mathbb{C}$ 上平凡 $\LG$-丛上的 Frenkel--Gross 联络：
\[
\nabla^{\LG}(\check X_0,\cdots,\check X_{r_\sigma})=d+\sum_{i=0}^{r_\sigma}\check X_i\,\frac{d}{dz},
\]
其中 $r_\sigma$ 是 $\mathbb{G}^\sigma$ 的秩，$\check X_i$ 是与 $\ckg$ 及 $\sigma$ 相伴的（扭）仿射 Kac--Moody 李代数的 $(-\alpha_i)$-根空间的一组基。

\begin{conjecture}
存在一般线性函数 $\varphi:I(1)/I(2)\to\mathbb{G}_{a,\mathbb{C}}$ 之集与基 $(\check X_0,\cdots,\check X_{r_\sigma})$ 之集之间的双射，使得只要 $\varphi$ 在这个双射下对应于 $(\check X_0,\cdots,\check X_{r_\sigma})$，就有 $\PP^1\setminus\{0,\infty\},\mathbb{C}$ 上 $\LG$-联络之间的自然同构
\[
\Kl^{\LG}(\varphi)_{dR}\cong({}^L G,\nabla^{\LG}(\check X_0,\cdots,\check X_{r_\sigma})).
\]
\end{conjecture}
